\section{Definitions and preliminary results}\label{sec:preliminaries}
\subsection{Norm conventions}
We use the unitary angular Fourier transform on $\R^3$ and Fourier series
$\widehat z(k)=\int_{\TT}z(x)e^{-2\pi i k\cdot x}\dd x$ on the unit torus.
Thus
\[
 \norm{z}_{H^s(\R^3)}^2=\int(1+|\xi|^2)^s|\widehat z(\xi)|^2\dd\xi,
 \qquad
 \norm{z}_{H^s(\TT)}^2=\sum_k(1+4\pi^2|k|^2)^s|\widehat z(k)|^2.
\]
These are the distributional Sobolev spaces of
Taylor~\cite[Chapters~3--4]{taylor2011}; vector and tensor norms sum the
squared component norms. The homogeneous norms replace these weights by
$|\xi|^{2s}$ and $(2\pi|k|)^{2s}$, respectively, with zero mean on the
torus. Tao's unnormalized Fourier convention gives the same Euclidean norms
after its Plancherel factor~\cite[Appendix~A]{taodispersive}.

Time norms are Bochner norms of strongly measurable paths, with equality
almost everywhere understood~\cite[Definitions~6.14, 6.27]{hunterpde}:
\begin{equation}\label{eq:time-norms}
 \norm{z}_{L^q(I;X)}=\left(\int_I\norm{z(t)}_X^q\dd t\right)^{1/q},
 \qquad \norm{z}_{L^\infty(I;X)}=\operatorname*{ess\,sup}_{t\in I}\norm{z(t)}_X.
\end{equation}
Force norms use $(0,\infty)$ unless stated otherwise, and give each smooth
force class its relative norm topology. On $(0,T)$ we write
\begin{equation}\label{eq:Enorm}
 \norm{z}_{E_T}=\norm{z}_{L^\infty_tL^2_x}+\norm{\nabla z}_{L^2_tL^2_x}.
\end{equation}
This is equivalent to the sum norm on $L^\infty_tL^2_x\cap L^2_tH^1_x$;
it assigns no value at the terminal time.

\subsection{Smooth data and solution classes}
All fields are real, and forces need not be divergence free. We write
$L^2_\sigma(D)$ for the divergence-free subspace of $L^2(D;\R^3)$.
The allowed initial velocities and forces are
\begin{align}
 \XX=\mathcal X_{\TT}&=C^\infty_{\mathrm{div}}(\TT;\R^3),&
 \FF=\mathcal F_{\TT}&=C_c^\infty(\TT\times(0,\infty);\R^3),
 \label{eq:inputspaces}\\
 \mathcal X_{\R}&=H^\infty(\R^3;\R^3)\cap L^2_\sigma(\R^3),
 \label{eq:Rinitial}
\end{align}
where $H^\infty=\bigcap_{m=0}^\infty H^m$ has its usual Fr\'echet topology.
On the whole space set
\begin{equation}\label{eq:Rclasses}
 \mathcal F_{\R}=\left\{f\in C^\infty([0,\infty);H^\infty):
 \norm{f}_{L^1_tH^m_x}+\norm{f}_{L^2_tH^m_x}<\infty
 \text{ for every integer }m\ge0\right\}.
\end{equation}
Smoothness into $H^\infty$ means smoothness into each $H^m$, with
one-sided time derivatives at zero. Thus periodic forces vanish near zero
and outside a compact time interval; whole-space forces may be nonzero at
zero and need not have compact support. Both classes consist of forces
defined through any possible blowup time of the velocity.

On the torus, velocity and pressure are periodic and $\int_{\TT}p=0$.
On $\R^3$, a classical velocity belongs to $C([0,S];H^m)$ for every
integer $m\ge0$ on each compact interval of its lifespan. Its pressure
gradient is specified below, and the scalar pressure is determined up to
a function of time. In either case local uniqueness defines a maximal
classical lifespan, denoted by $T^\nu_{\max}(a,f)$ on the torus and
$T^\nu_{\max,\R}(a,f)$ on the whole space. A given smooth solution is
\emph{smooth through $T$} if it extends smoothly to $[0,T+\delta]$
for some $\delta>0$ in the relevant class.

For fixed $\nu,T>0$, define
\begin{align}
 \BB_{\nu,a,T}&=\{f\in\FF:T^\nu_{\max}(a,f)\le T\},&
 \BB^0_{\nu,T}&=\BB_{\nu,0,T},\label{eq:singularforces}\\
 \mathcal B^\R_{\nu,a,T}
 &=\{f\in\mathcal F_{\R}:T^\nu_{\max,\R}(a,f)\le T\}.
 \label{eq:Rsingularforces}
\end{align}
We call these the forces producing classical breakdown by $T$, using
``breakdown for failure of classical continuation, as in the Clay
problem statement~\cite[alternatives (C) and (D)]{clay}.
The glued solutions satisfy the more specific condition of unbounded
speed at their terminal time.

\subsection{Homogeneous spaces and Fourier multipliers}
Write $\Lambda=(-\Delta)^{1/2}$ and $J=(I-\Delta)^{1/2}$, with the
Fourier conventions of Section~\ref{sec:intro}. On mean-zero periodic
fields, the homogeneous and inhomogeneous Sobolev norms are equivalent
at each fixed order. For $0<a<3/2$, the homogeneous completion is represented by
$\widehat z=|\xi|^{-a}G$, $G\in L^2$, with its unique
$L^{6/(3-2a)}$ representative supplied by Sobolev embedding.
The distributional pairing is well defined since $|\xi|^{-a}$ is locally
square integrable and Schwartz tests decay rapidly.
For the completed energy-force space we fix
\begin{equation}\label{eq:homogeneous-realization}
 \dot H^{-1}(\R^3)=\{h\in\mathcal S':\widehat h=F
 \text{ is a measurable function and }|\xi|^{-1}F\in L^2\}.
\end{equation}
Functions representing distributions are identified almost everywhere.
The maps $h\mapsto|\xi|^{-1}\widehat h$ and
$h\mapsto\langle\xi\rangle^s\widehat h$ identify these spaces with $L^2$;
the inverse Fourier data define tempered distributions. Thus they are
separable Hilbert spaces, with no polynomial ambiguity. Real vector fields
correspond to the conjugate-symmetric Fourier data.

\subsection{Standard Sobolev estimates}\label{sec:sobolev-estimates}
We use the tame product estimate and $H^2\hookrightarrow L^\infty$:
\begin{equation}\label{eq:Rproduct}
 \norm{vw}_{H^m}\le C_m(\norm{v}_{H^2}\norm{w}_{H^m}
              +\norm{w}_{H^2}\norm{v}_{H^m}),\qquad
 \norm{v}_\infty\le C\norm{v}_{H^2},\quad m\ge2.
\end{equation}
See Tao~\cite[Appendix~A, (A.13), Lemma~A.8]{taodispersive}; the periodic
version follows by the same Fourier convolution estimate, with
$\|\widehat v\|_{\ell^1}\le C\|v\|_{H^2}$. The critical embeddings used below are
\begin{equation}\label{eq:embeddings}
 \norm{v}_3\le C\norm{v}_{\dot H^{1/2}},\qquad
 \norm{\nabla v}_3+\norm{\Lambda v}_3\le C\norm{v}_{\dot H^{3/2}},\qquad
 \norm{\nabla v}_6\le C\norm{\Delta v}_2.
\end{equation}
They follow from the homogeneous Sobolev inequality
\cite[Appendix~A, (A.11)]{taodispersive}, applied also to derivatives.
On the torus use the compact-manifold embedding
\cite[Chapter~13, Proposition~6.4 and preceding discussion]{taylor2011iii}
and the spectral gap for mean-zero fields. These statements use the
homogeneous representatives fixed above; no embedding of $\dot H^{3/2}$
into $L^\infty$ is used.

\subsection{Projection, pressure, and local existence}\label{sec:analytic}
The Leray projection $\PP$ has Fourier symbol
$I-\xi\otimes\xi/|\xi|^2$ on $\R^3\setminus\{0\}$ and the same
symbol at $k\ne0$ on the torus. At the periodic zero mode it is the
identity. It is bounded on every $H^s$ and commutes with derivatives and
the heat semigroup. The projected equation is
\begin{equation}\label{eq:projected}
 \partial_tu-\nu\Delta u=-\PP\nabla\cdot(u\otimes u)+\PP f.
\end{equation}
On the torus, pressure is recovered from
\[
 \Delta p=\nabla\cdot f-\sum_{i,j}\partial_i\partial_j(u_i u_j),
 \qquad \int_{\TT}p=0.
\]
The right side has zero mean. On $\R^3$ we require
\begin{equation}\label{eq:Rpressure}
 \nabla p=(I-\PP)(f-\nabla\cdot(u\otimes u))=:G.
\end{equation}
The field $G$ is curl free. The radial potential
$p(x,t)=\int_0^1G(rx,t)\cdot x\dd r$ therefore recovers the pressure,
up to a function of time. We impose no $L^2$ condition on $p$ itself;
$\nabla p\in L^2$ excludes nonzero constant pressure gradients.

\begin{proposition}
\label{prop:local}\label{lem:Rlocal}
For each initial velocity in the stated class and each force smooth into
every $H^m$ on compact time intervals, equation~\eqref{eq:NS} has a unique
maximal smooth velocity, with pressure determined as above. If a solution
is defined on $[0,S)$, $S<\infty$, and
\begin{equation}\label{eq:criterion}
 \int_0^S\norm{u(t)}_{H^2(D)}^2\dd t<\infty,
\end{equation}
then it extends smoothly beyond $S$.
\end{proposition}
\begin{proof}
Apply Tao~\cite[Theorems~5.1(ii)--(iv), 5.4(ii)--(iv)]{tao2013} to
$\PP f$. The proof of Theorem~5.4(iv), p.~53, explicitly permits
$a\in H^k$ and $f\in C^j_tH^k_x$ for all $j,k$; the resulting smoothness
is on one common interval. Reduce viscosity to one by
$\tau=\nu t$, $\widetilde u=\nu^{-1}u$, and
$(\widetilde p,\widetilde f)=\nu^{-2}(p,f)$ with the corresponding time change.
For nonzero periodic mean, use Tao's mean reduction
\cite[Section~3, (36)--(38)]{tao2013}: set
$m(t)=\int_{\TT}a+\int_0^t\int_{\TT}f$ and translate by
$X(t)=\int_0^t m$, subtracting $m$ from the velocity and $m'$ from the force.
Pressure is recovered as above. Uniqueness and maximal development follow
from the same theorems and the overlap argument of Corollary~5.2.

For continuation, the standard $H^m$ energy estimate, Young's inequality
and \eqref{eq:Rproduct} give, for $m\ge3$,
\begin{equation}\label{eq:highcontinuation}
 (\norm{u}_{H^m})'\le C_{m,\nu}\norm{u}_{H^2}^2\norm{u}_{H^m}
                         +\norm{f}_{H^m}.
\end{equation}
Regularize the norm at zero. Gr\"onwall and \eqref{eq:criterion} bound
every $H^m$ norm uniformly up to $S$. Tao's quantitative local existence
then gives a uniform positive time when restarting at $t_0\uparrow S$;
the fixed force is smooth on $[0,S+1]$. Uniqueness glues one such solution
past $S$.
\end{proof}

\subsection{Energy and initial vanishing of the compact solution}
We use the fixed solution $(U,P,F)$ chosen after
Theorem~\ref{thm:packet}. Its energy estimate supplies the time-integrated
dissipation bound and the initial vanishing needed for rescaling.

\begin{lemma}\label{lem:packetenergy}
The building block satisfies
\[
 M:=\sup_{0\le t<1}\norm{U(t)}_2<\infty,\qquad
 D:=\norm{\nabla U}_{L^2((0,1)\times\R^3)}<\infty.
\]
Both $U$ and $P$ vanish on an initial time interval and therefore extend smoothly by zero to negative times.
\end{lemma}
\begin{proof}
The classical energy identity and $U(0)=0$ give, with
$N(t)=\int_0^t\|F(s)\|_2\dd s$, the bounds $\|U(t)\|_2\le N(t)$ and
\begin{equation}\label{eq:packetenergy}
 \norm{U(t)}_2^2+2\nu\int_0^t\norm{\nabla U(s)}_2^2\dd s
 \le 2\int_0^t\norm{F(s)}_2N(s)\dd s=N(t)^2.
\end{equation}
Compact support justifies the identity; regularizing $\|U\|_2$ justifies
the first bound at zero. Since $F$ is time compact, monotone convergence
gives finite dissipation up to time one. Before the support of $F$, the
same estimate gives $U=0$; the equation and compact support then give $P=0$.
Both fields therefore extend smoothly by zero to negative times.
\end{proof}
